% Layered architecture with an explicit trust boundary. The proven
% kernel is the bottom, smallest layer; expressiveness grows upward and
% sits outside the trusted base.
\begin{figure}[t]
\centering
\begin{tikzpicture}[
    font=\footnotesize,
    layer/.style={draw, rounded corners, minimum width=88mm, align=center,
                  inner sep=4pt, minimum height=11mm},
    trusted/.style={layer, fill=green!12},
    declarative/.style={layer, fill=black!6},
    learned/.style={layer, fill=blue!9},
    >={Stealth[length=2mm]},
    node distance=3.5mm,
]
\node[learned]                       (llm)  {\textbf{Authoring (learned, untrusted)}\\
                     LLM intent capture $\to$ candidate specifications};
\node[declarative, below=of llm]     (cp)   {\textbf{Control plane (declarative)}\\
                     content-addressed posting rules, mappings, policy};
\node[declarative, below=of cp]      (facts){\textbf{Fact plane (declarative)}\\
                     typed prices / FX with provenance + validity windows};
\node[trusted, below=8mm of facts]   (ker)  {\textbf{Conservation kernel (proven, Lean~4)}\\[1pt]
                     $\rconserves(\rvec):\rvec=\rzero$ \quad
                     $\rledger=\mathrm{foldl}\,(+)\,\rzero\,\textit{log}$};

% Trust boundary between the declarative planes and the proven kernel.
\draw[dashed, thick, red!60!black]
    ([yshift=4mm]ker.north west) -- ([yshift=4mm]ker.north east)
    node[midway, above, font=\scriptsize\itshape, text=red!60!black]
    {trust boundary --- only conserving transactions cross downward};

\draw[->] (llm) -- (cp);
\draw[->] (cp) -- (facts);
\draw[->] (facts) -- (ker);
\end{tikzpicture}
\caption{The trust boundary. Expressiveness increases upward; the trusted
computing base is only the bottom layer --- the conservation kernel and
the theorems over it. Learned and declarative layers above produce
\emph{proposals}; the kernel admits a transaction only if it conserves
value, so nothing above the boundary can unbalance the ledger.}
\label{fig:layers}
\end{figure}
